% Pengesahan Tesis (Lampiran 12 - Panduan UDINUS)
\thispagestyle{plain}
\begin{center}
\addcontentsline{toc}{chapter}{PENGESAHAN TESIS}
% Logo di atas
\IfFileExists{\UniversityLogo}{\includegraphics[width=2.8cm]{\UniversityLogo}\\[0.4cm]}{\relax}
{\bfseries \UniversityName}\\[0.6cm]
{\bfseries\large PENGESAHAN TESIS}\\[0.6cm]

\begin{tabular}{@{}l@{\hspace{0.5cm}}c@{\hspace{0.5cm}}p{10cm}@{}}
JUDUL & : & {\bfseries \TitleID} \\[0.2cm]
NAMA  & : & \AuthorNameTitleCase \\[0.2cm]
NIM   & : & \StudentID \\
\end{tabular}

\vspace{0.5cm}
Tesis ini telah diujikan dan dipertahankan di hadapan Dewan Penguji pada
Sidang Tesis tanggal \ExamDate.
Menurut pandangan kami, Tesis ini
memadai dari segi kualitas untuk tujuan penganugerahan gelar
{\DegreeName} (M.Kom.)\\[0.5cm]

{\ApprovalPlace},\quad\RatificationDate\\[1.5cm]

Dewan Penguji:\\[1.4cm]

\begin{tabular*}{0.86\textwidth}{@{\extracolsep{\fill}}cc@{}}
    % Lebar garis tanda tangan diatur oleh \rule{5.5cm}{0.4pt}.
    % Ubah 5.5cm untuk memperpanjang/memendekkan garis.
    \makebox[5.5cm][c]{\rule{5.5cm}{0.4pt}} & \makebox[5.5cm][c]{\rule{5.5cm}{0.4pt}} \\
    \makebox[5.5cm][c]{\ExaminerA} & \makebox[5.5cm][c]{\ExaminerB} \\
    \makebox[5.5cm][c]{Anggota} & \makebox[5.5cm][c]{Anggota} \\[1.2cm]
\end{tabular*}

\begin{tabular*}{0.86\textwidth}{@{\extracolsep{\fill}}cc@{}}
    % "Mengetahui" ditempatkan di kolom kiri (di atas Dekan) sesuai contoh
    % Lampiran 12 panduan UDINUS, bukan di tengah halaman.
    \makebox[5.5cm][c]{Mengetahui} & \\[1.4cm]
    \makebox[5.5cm][c]{\rule{5.5cm}{0.4pt}} & \makebox[5.5cm][c]{\rule{5.5cm}{0.4pt}} \\
    \makebox[5.5cm][c]{\DeanName} & \makebox[5.5cm][c]{\ChiefExaminer} \\
    \makebox[5.5cm][c]{Dekan} & \makebox[5.5cm][c]{Ketua Penguji} \\
\end{tabular*}
\end{center}
\clearpage